% !TEX program = xelatex
\documentclass[12pt, a4paper]{article}

% === 中文支持 ===
\usepackage{xeCJK}
\setCJKmainfont{SimSun}          % 宋体
\setCJKsansfont{SimHei}          % 黑体
\setCJKmonofont{FangSong}        % 仿宋

% === 基础包 ===
\usepackage[utf8]{inputenc}
\usepackage[T1]{fontenc}
\usepackage{geometry}
\usepackage{setspace}
\usepackage{amsmath}
\usepackage{graphicx}
\usepackage{booktabs}
\usepackage{hyperref}
\usepackage{url}
\usepackage{float}
\usepackage{longtable}
\usepackage{array}
\usepackage{fancyhdr}
\usepackage{enumitem}
\usepackage{caption}

% === 页面设置（A4纸，上下左右各2.5厘米）===
\geometry{
    a4paper,
    top=2.5cm,
    bottom=2.5cm,
    left=2.5cm,
    right=2.5cm,
    headheight=14pt
}

% === 行距设置 ===
\onehalfspacing
\setlength{\parindent}{2em}

% === 超链接设置 ===
\hypersetup{
    colorlinks=true,
    linkcolor=black,
    citecolor=black,
    urlcolor=blue,
    breaklinks=true
}

% === 页眉页脚设置 ===
\pagestyle{fancy}
\fancyhf{}
\fancyfoot[C]{\thepage}
\renewcommand{\headrulewidth}{0pt}
\renewcommand{\footrulewidth}{0pt}

% === 标题格式设置 ===
\usepackage{titlesec}
\titleformat{\section}{\centering\zihao{4}\heiti}{\thesection}{1em}{}
\titleformat{\subsection}{\zihao{-4}\heiti}{\thesubsection}{1em}{}
\titleformat{\subsubsection}{\zihao{5}\heiti}{\thesubsubsection}{1em}{}

% === 字号设置命令 ===
\newcommand{\zihao}[1]{%
    \ifnum#1=0 {\fontsize{42pt}{50pt}\selectfont}%
    \else\ifnum#1=-0 {\fontsize{36pt}{44pt}\selectfont}%
    \else\ifnum#1=1 {\fontsize{26pt}{32pt}\selectfont}%
    \else\ifnum#1=-1 {\fontsize{24pt}{28pt}\selectfont}%
    \else\ifnum#1=2 {\fontsize{22pt}{26pt}\selectfont}%
    \else\ifnum#1=-2 {\fontsize{18pt}{22pt}\selectfont}%
    \else\ifnum#1=3 {\fontsize{16pt}{20pt}\selectfont}%
    \else\ifnum#1=-3 {\fontsize{15pt}{18pt}\selectfont}%
    \else\ifnum#1=4 {\fontsize{14pt}{17pt}\selectfont}%
    \else\ifnum#1=-4 {\fontsize{12pt}{15pt}\selectfont}%
    \else\ifnum#1=5 {\fontsize{10.5pt}{13pt}\selectfont}%
    \else\ifnum#1=-5 {\fontsize{9pt}{11pt}\selectfont}%
    \else\ifnum#1=6 {\fontsize{7.5pt}{9pt}\selectfont}%
    \else\ifnum#1=-6 {\fontsize{6.5pt}{8pt}\selectfont}%
    \else\ifnum#1=7 {\fontsize{5.5pt}{7pt}\selectfont}%
    \else\ifnum#1=-7 {\fontsize{5pt}{6pt}\selectfont}%
    \fi\fi\fi\fi\fi\fi\fi\fi\fi\fi\fi\fi\fi\fi\fi\fi
}

\newcommand{\heiti}{\CJKfontspec{SimHei}}
\newcommand{\songti}{\CJKfontspec{SimSun}}
\newcommand{\kaishu}{\CJKfontspec{KaiTi}}
\newcommand{\fangsong}{\CJKfontspec{FangSong}}

% === 元数据 ===
\title{\zihao{-2}\heiti 微小卫星遥测智能异常检测与RAG解释辅助系统}
\author{\zihao{4}\kaishu 李冰杰$^{1}$ \quad 王翌钧$^{2}$ \quad 陶诗怡$^{3}$ \\[6pt]
  \zihao{5}\songti $^{123}$上海电机学院 电子信息学院，上海 201306}
\date{}

% === 文档开始 ===
\begin{document}

\maketitle

% === 摘要 ===
\begin{center}
\zihao{5}\heiti 摘\quad 要
\end{center}

\zihao{5}\songti
微小卫星在轨运行期间产生海量遥测数据，人工排查异常耗时且易遗漏。传统方法依赖固定阈值和专家规则，难以捕捉多维遥测参数间的非线性耦合关系。本研究提出分通道自适应异常检测框架，采用``检测器—策略''解耦的模块化设计，针对磁力计强通道和光电二极管弱通道的特性差异，分别采用Isolation Forest与统计规则兜底的差异化检测逻辑。基于ESA OPS-SAT-AD公开数据集（2,123个遥测段、9个通道、20.4\%异常率）的系统性实验表明：(1) 分通道建模将SegF1从全局基线的0.429提升至0.545（+27.0\%）；(2) 规则兜底融合进一步提升至0.576（+5.7\%）；(3) 阈值优化达到最终0.578（+0.3\%），整体相对提升约34.6\%。消融实验验证了18维段级特征优于1维均值特征（+68\%--128\%）。特征重要性分析通过IForest内置和置换重要性双方法对比，揭示了峰值特征（smooth10/20\_n\_peaks）对异常检测的关键作用。结合检索增强生成（RAG）技术构建异常解释系统，知识库包含5,062个文档块，检索平均Top-1相似度0.640，生成成功率100\%，可为异常检测结果生成自然语言形式的专家级诊断报告。

\vspace{6pt}
\noindent
\zihao{5}\heiti 关键词\quad \zihao{5}\songti 微小卫星；异常检测；Isolation Forest；分通道建模；检索增强生成

\newpage
\setcounter{page}{1}

% === 一、引言 ===
\section{引言}

\subsection{研究背景}

\zihao{5}\songti
随着航天技术的快速发展，微小卫星（Microsatellite）因其成本低、研制周期短、发射灵活等优势，在遥感、通信、科学探测等领域得到广泛应用[1]。微小卫星在轨运行期间，星上各子系统持续产生海量遥测数据，这些数据反映了卫星的运行状态和健康状况。然而，人工排查遥测数据中的异常模式耗时且易遗漏，尤其是在数据量庞大、参数维度高的情况下，传统的人工分析方法已难以满足实时性和准确性的要求。

遥测异常检测是卫星健康管理的核心环节，其目标是从正常运行数据中识别出偏离预期行为的异常模式。早期的异常检测方法主要依赖人工规则和专家知识，如固定阈值法（$\sigma$规则）、包络线法、符号化专家系统等。这些方法虽然实现简单、可解释性强，但存在以下根本局限：

\begin{enumerate}[label=(\arabic*)]
    \item \textbf{单参数阈值限制}：普遍依赖单参数阈值判断，难以捕捉多维遥测参数之间的非线性耦合关系。例如，磁力计三轴数据之间存在复杂的几何关系，单参数阈值无法有效建模这种耦合。
    \item \textbf{规则完备性问题}：专家系统虽能编码部分耦合规则，但面对海量参数组合，规则完备性难以保证。随着参数维度增加，规则数量呈指数增长，维护成本极高。
    \item \textbf{静态判据局限}：静态阈值和规则在动态工况下虚警率较高。卫星在不同轨道位置、不同工作模式下，遥测数据的统计特性会发生显著变化，静态判据难以适应这种动态性。
\end{enumerate}

为克服上述局限，研究者引入机器学习方法进行遥测异常检测。传统的无监督异常检测方法包括One-Class SVM（OCSVM）[3]、基于LSTM的编码器-解码器架构[4]等。OCSVM通过构建正常数据的边界来识别异常，但在高维数据上计算复杂度较高；LSTM编码器-解码器通过时序建模捕获传感器间的依赖关系，但训练成本高且可解释性差。近年来，时序卷积网络（TCN）[5]在序列建模任务中展现出与循环网络相当的性能，且具有更好的并行性。

\subsection{问题提出}

\zihao{5}\songti
本项目采用无监督学习算法对欧空局（ESA）OPS-SAT卫星遥测数据进行异常检测。在实际分析中发现，微小卫星遥测参数间存在极强的异构性：磁力计通道呈现规律性的周期信号，而光电二极管通道则表现为高噪声的脉冲信号。这种异构性导致``一刀切''的全局模型往往顾此失彼——对磁力计通道有效的特征提取策略，在光电二极管通道上可能完全失效。

针对这一挑战，本研究提出分通道自适应检测框架，通过``检测器-策略''解耦设计，为不同特性的通道定制差异化检测逻辑与规则兜底。同时，为提升异常检测结果的可解释性，本研究结合检索增强生成（Retrieval-Augmented Generation, RAG）技术构建异常解释系统，为运维人员提供自然语言形式的异常原因分析。

\subsection{研究目标与贡献}

\zihao{5}\songti
本研究的核心目标是：在保持框架通用性的前提下，针对不同遥测通道的数据特性差异，设计有效的分通道自适应异常检测策略，并构建可解释的异常解释系统。

本研究的主要贡献包括：

\begin{enumerate}[label=(\arabic*)]
    \item 提出分通道自适应异常检测框架，采用``检测器-策略''解耦的模块化设计，支持针对不同通道特性定制差异化检测逻辑。
    \item 设计强弱通道分离策略，对磁力计强通道采用IForest+规则OR融合，对光电二极管弱通道采用纯规则兜底，有效应对数据异构性挑战。
    \item 构建基于RAG技术的异常解释系统，知识库包含5,062个文档块，可为异常检测结果生成自然语言解释。
    \item 基于ESA OPS-SAT-AD公开数据集进行系统性实验验证，SegF1从全局基线的0.429提升至0.578，相对提升34.6\%。
\end{enumerate}

% === 二、相关工作 ===
\section{相关工作}

\subsection{卫星遥测异常检测}

\zihao{5}\songti
卫星遥测异常检测是航天器健康管理的关键技术，已有数十年的研究历史。早期方法主要基于统计学和信号处理技术，如CUSUM（累积和）控制图、EWMA（指数加权移动平均）等。这些方法假设数据服从特定分布，对非高斯、非平稳数据的适应性较差。

近年来，机器学习方法在卫星遥测异常检测中得到广泛应用。根据学习方式的不同，可分为监督学习、无监督学习和半监督学习三类：

\textbf{监督学习方法}：需要标注的异常样本进行训练。AdaBoost[13]、随机森林等集成方法在OPS-SAT基准测试中取得了较高的准确率（F1达0.836），但依赖大量标注数据，而实际场景中异常样本往往稀缺[1]。

\textbf{无监督学习方法}：仅使用正常数据训练模型，无需异常标注。Isolation Forest[2]通过随机分割隔离异常点，具有线性时间复杂度；OCSVM[3]构建正常数据的超球面边界；局部离群因子（LOF）[8]基于密度比识别局部异常；自编码器（Autoencoder）通过重构误差识别异常。Ruff等人[18]和Pang等人[19]对深度与浅层异常检测方法进行了全面综述，指出无监督方法在标注数据稀缺的场景下具有独特优势。

\textbf{半监督学习方法}：结合少量标注数据和大量无标注数据进行训练。这类方法在实际应用中具有较好的平衡性，但实现复杂度较高。

在卫星遥测领域，Hundman等人[7]提出了基于LSTM的航天器异常检测方法，在NASA SMAP和MSL数据集上取得了显著效果，其非参数动态阈值策略对本研究的阈值优化设计提供了重要启发。Castellani等人[17]探索了数字孪生技术在卫星异常检测中的应用。Ruszczak等人[1]发布的OPS-SAT-AD数据集为本研究提供了标准化的基准测试平台。

\subsection{无监督异常检测算法}

\zihao{5}\songti
无监督异常检测算法可分为基于统计的方法、基于距离的方法、基于密度的方法和基于树的方法四大类[18]。

\textbf{基于统计的方法}：假设数据服从特定分布（如正态分布），将低概率区域的数据点判定为异常。3$\sigma$规则和IQR规则是这类方法的典型代表。优点是计算简单、可解释性强，缺点是对非正态分布数据的适应性较差。

\textbf{基于距离的方法}：通过计算数据点之间的距离来识别异常。k近邻（kNN）算法认为远离其k个近邻的数据点为异常。这类方法的计算复杂度为$O(n^2)$，在大规模数据上效率较低。

\textbf{基于密度的方法}：LOF[8]通过比较数据点的局部密度与其邻域密度来识别异常。LOF能够检测局部异常（即在某些区域密度显著低于邻域的点），但对参数$k$的选择敏感。

\textbf{基于树的方法}：Isolation Forest[2]是这类方法的代表。与上述方法不同，IForest不依赖距离或密度计算，而是通过随机分割来隔离异常点，具有$O(tn)$的时间复杂度（$t$为树的数量，$n$为样本数，当采样大小$\psi$固定时简化为$O(n)$）和$O(n)$的空间复杂度，适合大规模数据场景。

在OPS-SAT-AD基准测试中，各算法的表现差异显著[1]：IForest的SegF1为0.295，OCSVM为0.647，MO-GAAL为0.726，AdaBoost（监督）为0.836。这种差异反映了不同算法对数据特性的适应性不同，也为本研究的分通道策略提供了算法选择的参考依据。

\subsection{时序异常检测}

\zihao{5}\songti
遥测数据本质上是时间序列，因此时序异常检测方法具有天然的适用性。

\textbf{基于循环网络的方法}：Malhotra等人[4]提出了基于LSTM编码器-解码器的多传感器异常检测框架，通过学习正常时序模式的低维表示来识别异常。该方法能够捕获传感器间的时序依赖关系，但训练成本高且可解释性差。

\textbf{基于卷积的方法}：Bai等人[5]通过实验证明时序卷积网络（TCN）在序列建模任务中展现出与循环网络相当的性能，且具有更好的并行性和更长的记忆能力。TCN通过因果卷积和膨胀卷积实现对长序列的高效建模。

\textbf{基于Transformer的方法}：近年来，Transformer架构在时序异常检测中也得到应用。通过自注意力机制，Transformer能够捕获序列中的长距离依赖关系，但计算复杂度较高。

本研究采用段级特征提取策略，将时序数据转化为统计特征向量，从而避免了直接处理变长时间序列的复杂性，同时保留了关键的时序信息（如峰值数、差分特征等）。

\subsection{Isolation Forest算法}

\zihao{5}\songti
Isolation Forest（IF）[2]是一种基于树结构的无监督异常检测算法，由Liu等人于2008年提出。其核心思想是：异常点由于与正常数据点在特征空间中的分布不同，更容易被随机分割所隔离。

算法的基本流程如下：

\begin{enumerate}[label=(\arabic*)]
    \item \textbf{构建隔离树}：从训练数据中随机选择一个特征和分割值，递归地将数据分割成两个子集，直到每个子集只包含一个样本或达到最大深度。
    \item \textbf{计算路径长度}：对于每个样本，计算其在隔离树中从根节点到叶子节点的路径长度。异常点由于分布稀疏，路径长度通常较短。
    \item \textbf{计算异常分数}：综合多棵隔离树的结果，计算每个样本的异常分数。异常分数越高，样本越可能是异常。
\end{enumerate}

IF的异常分数定义为：

\begin{equation}
s(x, n) = 2^{-\frac{E[h(x)]}{c(n)}}
\end{equation}

其中，$h(x)$为样本$x$在隔离树中的路径长度，$E[h(x)]$为多棵树的平均路径长度，$c(n) = 2H(n-1) - 2(n-1)/n$为归一化因子（$H(i)$为调和数），$n$为训练样本数。当$s \to 1$时表示高度异常，$s \to 0.5$时表示正常，$s \ll 0.5$时表示被隔离在密集区域。

IF具有以下优势：
\begin{itemize}
    \item 线性时间复杂度$O(n\log n)$，适合大规模数据
    \item 对高维数据适应性好
    \item 内存占用低，适合资源受限环境
    \item 实现简单，易于部署
    \item 无需距离或密度计算，避免了维度灾难
\end{itemize}

\subsection{统计异常检测规则}

\zihao{5}\songti
除了机器学习方法，统计规则在异常检测中也发挥着重要作用。常用的统计规则包括：

\textbf{3$\sigma$规则}：假设数据服从正态分布，将超出均值$\pm 3$倍标准差范围的数据点判定为异常。该规则简单直观，但对非正态分布数据的适应性较差。在本研究中，我们将阈值设为$\pm 3\sigma$，即：
\begin{equation}
\text{异常} \Leftrightarrow |x_i - \bar{x}| > 3\sigma
\end{equation}

\textbf{IQR规则}：基于四分位距（Interquartile Range）的异常检测方法。将超出$Q1 - 1.5 \times IQR$或$Q3 + 1.5 \times IQR$范围的数据点判定为异常，其中$IQR = Q3 - Q1$。该规则对异常值具有更好的鲁棒性，因为四分位数对极端值不敏感。

在本研究中，我们将3$\sigma$规则和IQR规则组合作为统计兜底策略，对每个通道的每个特征独立检查，各项触发独立计数。单个通道的累计违反次数$\geq 2$即判定该段为异常。通过逻辑``OR''集成多个规则的判断结果，有效提升召回率。

\subsection{检索增强生成（RAG）技术}

\zihao{5}\songti
检索增强生成（Retrieval-Augmented Generation, RAG）是Lewis等人[10]于2020年提出的一种结合信息检索和文本生成的技术范式。RAG系统通过从外部知识库中检索相关文档，将其作为上下文提供给大语言模型（LLM），从而生成更准确、更有依据的回答。

RAG系统的典型架构包括三个核心组件：

\begin{enumerate}[label=(\arabic*)]
    \item \textbf{文档索引}：将外部知识文档分割成块（chunks），通过嵌入模型转换为向量表示，存储在向量数据库中。本研究采用BGE-M3[11]作为嵌入模型，FAISS[12]作为向量检索引擎。
    \item \textbf{检索模块}：根据用户查询，从向量数据库中检索语义最相关的文档块。相似度度量通常采用余弦相似度或内积。
    \item \textbf{生成模块}：将检索到的文档块与用户查询一起输入大语言模型，生成最终回答。本研究使用DeepSeek-V3作为生成模型。
\end{enumerate}

RAG技术的优势在于：
\begin{itemize}
    \item 减少大语言模型的幻觉（hallucination）问题
    \item 支持知识库的动态更新
    \item 提供可追溯的信息来源
    \item 降低模型训练成本
\end{itemize}

在本研究中，我们将RAG技术应用于异常解释生成，通过检索相关的技术文档和历史案例，为异常检测结果提供自然语言形式的解释。知识库包含5,062个文档块，涵盖SCI论文、NASA技术报告、ITU标准和实验知识文档等多种来源。

% === 三、数据集 ===
\section{数据集}

\subsection{数据集概述}

\zihao{5}\songti
本研究使用ESA OPS-SAT-AD公开数据集[1]，该数据集由欧洲航天局（ESA）发布，旨在为卫星遥测异常检测研究提供标准化的基准测试平台。OPS-SAT是一颗3单位的立方星，于2019年发射升空，主要用于在轨技术验证。

数据集的数据管线统一使用segments.csv，共包含2,123个遥测数据段（segments），覆盖9个遥测通道。每个数据段标注为正常（anomaly=0）或异常（anomaly=1），异常样本占比约20.4\%，呈明显的类别不平衡特征。

\subsection{遥测通道分类}

\zihao{5}\songti
9个遥测通道均来自卫星有效载荷子系统，按数据特性划分为强通道和弱通道两类，如表~\ref{tab:channels}所示。

\begin{table}[H]
\centering
\caption{数据集遥测通道分类}
\label{tab:channels}
\begin{tabular}{llll}
\toprule
通道ID & 类型 & 物理含义 & 测试段数 \\
\midrule
CADC0872 & 强 & 磁力计X轴 & 132 \\
CADC0873 & 强 & 磁力计Y轴 & 153 \\
CADC0874 & 强 & 磁力计Z轴 & 52 \\
CADC0884 & 弱 & 光电二极管1 & - \\
CADC0886 & 弱 & 光电二极管2 & 4 \\
CADC0888 & 弱 & 光电二极管3 & 64 \\
CADC0890 & 弱 & 光电二极管4 & 2 \\
CADC0892 & 弱 & 光电二极管5 & 53 \\
CADC0894 & 弱 & 光电二极管6 & 33 \\
\bottomrule
\end{tabular}
\end{table}

\subsection{强弱通道特性分析}

\zihao{5}\songti
按数据动力学特性将9个通道划分为两类：

\textbf{高信噪比通道（强特征）}：以磁力计三轴（CADC0872/0873/0874）为代表，信号具有明显的物理规律。磁力计数据反映卫星在地磁场中的姿态变化，呈现周期性的正弦波形，信噪比较高，适合无监督模型提取特征。强通道的测试段数较多（132、153、52段），异常样本相对充足，为模型训练提供了充分的数据支持。

\textbf{低信噪比/非平稳通道（弱特征）}：以光电二极管类（CADC0884-CADC0894）为代表，受外部环境干扰大，噪声极高且数据分布稀疏。光电二极管用于测量太阳光强度，受卫星姿态、阴影遮挡、太阳角度等多种因素影响，信号呈现高度非平稳性。弱通道的测试段数普遍较少（2-64段），异常样本极度匮乏（如CADC0890仅有2段），对纯机器学习模型挑战巨大。

\subsection{数据分布特征}

\zihao{5}\songti
如表~\ref{tab:channel_stats}所示，强弱通道在数据分布上呈现极强的异构性。

\begin{table}[H]
\centering
\caption{各通道测试集统计特征}
\label{tab:channel_stats}
\begin{tabular}{lcccc}
\toprule
通道 & 测试段数 & 异常段数 & 异常占比 & 通道类型 \\
\midrule
CADC0872 & 132 & 32 & 24.2\% & 强 \\
CADC0873 & 153 & 31 & 20.3\% & 强 \\
CADC0874 & 52 & 23 & 44.2\% & 强 \\
CADC0886 & 4 & 1 & 25.0\% & 弱 \\
CADC0888 & 64 & 12 & 18.8\% & 弱 \\
CADC0890 & 2 & 2 & 100.0\% & 弱 \\
CADC0892 & 53 & 7 & 13.2\% & 弱 \\
CADC0894 & 33 & 5 & 15.2\% & 弱 \\
\bottomrule
\end{tabular}
\end{table}

特别是光电二极管类弱通道（如CADC0886、CADC0890），不仅测试段总数稀疏，异常段落更是极度匮乏（分别仅有1段和2段）。这种极端的数据分布为后续``在弱通道弃用无监督模型、直接采用统计规则兜底''的策略提供了直接的先验依据，避免了模型在小样本噪声中的过拟合风险。

% === 四、方法 ===
\section{方法}

\subsection{框架总体架构}

\zihao{5}\songti
本研究提出的分通道自适应异常检测框架采用``检测器—策略''解耦的模块化设计，其总体架构如图~\ref{fig:framework}所示。

\begin{figure}[H]
\centering
\framebox[0.9\textwidth]{框架架构图占位}
\caption{分通道自适应异常检测框架架构}
\label{fig:framework}
\end{figure}

框架的核心设计思想是将异常检测过程分解为三个层次：

\begin{enumerate}[label=(\arabic*)]
    \item \textbf{特征提取层}：从原始遥测数据中提取18维段级特征，包括统计特征、峰值特征、差分特征等。
    \item \textbf{检测器层}：核心异常检测器具有算法无关性，可灵活替换。本研究采用Isolation Forest作为默认检测器。
    \item \textbf{策略层}：根据通道特性选择差异化的检测策略，包括纯模型检测、模型+规则融合、纯规则检测等。
\end{enumerate}

这种解耦设计的核心优势在于：当需要升级检测器算法时，只需替换检测器层，而无需修改策略层和特征提取层，从而保证了框架的通用性和可扩展性。

\subsection{特征工程}

\zihao{5}\songti
本研究从原始遥测数据中提取18维段级特征，这些特征从不同角度描述了数据段的统计特性和形态特征。特征定义如表~\ref{tab:features}所示。

\begin{table}[H]
\centering
\caption{18维段级特征定义}
\label{tab:features}
\begin{tabular}{lll}
\toprule
特征名称 & 类型 & 描述 \\
\midrule
duration & 基础 & 数据段时长 \\
len & 基础 & 数据点数量 \\
mean & 统计 & 均值 \\
var & 统计 & 方差 \\
std & 统计 & 标准差 \\
kurtosis & 统计 & 峰度 \\
skew & 统计 & 偏度 \\
n\_peaks & 峰值 & 原始信号峰值数 \\
smooth10\_n\_peaks & 峰值 & 10点平滑后峰值数 \\
smooth20\_n\_peaks & 峰值 & 20点平滑后峰值数 \\
diff\_peaks & 差分 & 一阶差分峰值数 \\
diff2\_peaks & 差分 & 二阶差分峰值数 \\
diff\_var & 差分 & 一阶差分方差 \\
diff2\_var & 差分 & 二阶差分方差 \\
gaps\_squared & 间隔 & 间隔平方和 \\
len\_weighted & 加权 & 加权长度 \\
var\_div\_duration & 比值 & 方差/时长 \\
var\_div\_len & 比值 & 方差/数据点数 \\
\bottomrule
\end{tabular}
\end{table}

\subsection{分通道检测策略}

\zihao{5}\songti
针对各通道的数据特性差异，本框架将9个有效遥测通道划分为强、弱两类，并采取差异化的检测逻辑：

\subsubsection{强通道策略}

\textbf{强通道（CADC0872/0873/0874，磁力计三轴）}：采用IF + 统计规则OR融合的策略。

IForest检测器配置如下：
\begin{itemize}
    \item contamination参数：0.1（预估异常比例）
    \item n\_estimators：100（隔离树数量）
    \item max\_samples：256（每棵树的采样数）
    \item random\_state：42（随机种子，保证可复现）
\end{itemize}

统计规则的具体实现为：对每个通道的每个特征，分别检查其值是否超出该通道训练集的$3\sigma$界限和IQR界限（$Q1-1.5\times IQR$, $Q3+1.5\times IQR$），各项触发独立计数。单个通道的累计违反次数$\geq 2$即判定该段为异常。

两者通过逻辑``OR''集成，任一方法报警即判为异常。这种融合策略的优势在于：
\begin{itemize}
    \item IForest负责捕获多维非线性异常模式
    \item 统计规则负责捕获单维度的极端异常
    \item OR融合有效提升召回率，同时控制虚警在可接受范围
\end{itemize}

\subsubsection{弱通道策略}

\textbf{弱通道（CADC0884/0886/0888/0890/0892/0894，光电二极管类）}：IForest在该类通道上因数据噪声极大、规律性弱而导致效果接近随机猜测，因此仅使用规则兜底（$3\sigma$ + IQR）作为主判断。

弱通道策略的设计依据在于：
\begin{itemize}
    \item 数据规律性极弱，IForest容易在噪声中产生虚假关联
    \item 测试段数稀疏，模型训练数据不足
    \item 异常样本极度匮乏，模型难以学习有效的异常模式
    \item 统计规则具有更好的鲁棒性和可解释性
\end{itemize}

这种``降级处理''思路，通过鲁棒性更强的统计规则进行拦截，旨在以极低的计算代价换取更可靠的召回，避免复杂模型在弱特征数据上的``过拟合''倾向。

\subsection{阈值优化}

\zihao{5}\songti
在完成分通道检测后，需要对各通道的异常分数进行融合。本研究采用阈值优化策略，通过调整异常判定阈值来平衡精确率和召回率。

阈值优化的过程如下：
\begin{enumerate}[label=(\arabic*)]
    \item 在$[0.1, 0.5]$范围内以0.05为步长搜索最优阈值
    \item 对每个阈值，计算融合后的SegF1
    \item 选择SegF1最大的阈值作为最终阈值
\end{enumerate}

实验表明，最优阈值为0.1，此时SegF1达到0.578。

\subsection{RAG解释系统架构}

\zihao{5}\songti
本研究构建的RAG异常解释系统架构如图~\ref{fig:rag}所示。

\begin{figure}[H]
\centering
\framebox[0.9\textwidth]{RAG系统架构图占位}
\caption{RAG异常解释系统架构}
\label{fig:rag}
\end{figure}

系统包含以下核心组件：

\subsubsection{知识库构建}

知识库包含5,062个文档块（chunks），来源包括：
\begin{itemize}
    \item Ruszczak等人[1]的SCI论文：提供数据集定义和基准测试方法
    \item NASA SOA 2024报告：提供卫星异常检测的最佳实践
    \item MAXWELL 2023手册：提供卫星遥测数据处理规范
    \item ITU-R报告：提供无线电通信相关标准
    \item 实验知识MD文档：记录项目实验过程和经验总结
\end{itemize}

文档分割采用固定大小分块策略，块大小为512 tokens，重叠64 tokens，确保上下文连续性。

\subsubsection{向量索引}

使用嵌入模型将文档块转换为向量表示，存储在向量数据库中。嵌入模型选择考虑了以下因素：
\begin{itemize}
    \item 支持中英文双语
    \item 向量维度可配置（本研究使用768维，通过Matryoshka表示学习从默认1024维截断）
    \item 推理速度快
    \item 内存占用低
\end{itemize}

\subsubsection{检索与生成}

当异常检测系统标记异常段后，系统执行以下步骤：
\begin{enumerate}[label=(\arabic*)]
    \item 提取异常段的特征描述和上下文信息
    \item 将描述转换为向量，在知识库中检索最相关的k个文档块
    \item 将检索到的文档块与异常信息一起输入大语言模型
    \item 大语言模型生成自然语言形式的异常解释
\end{enumerate}

% === 五、实验 ===
\section{实验}

\subsection{实验设置}

\subsubsection{数据集划分}

\zihao{5}\songti
所有实验均严格遵循OPS-SAT-AD数据集中预设的train列进行训练/测试集划分，保证与基准发布者的评估协议一致。训练集用于模型训练和参数调优，测试集用于性能评估。

\subsubsection{评估指标}

本研究采用以下评估指标：

\textbf{段级F1（SegF1）}：主要评估指标，计算公式为：
\begin{equation}
SegF1 = \frac{2 \times Precision \times Recall}{Precision + Recall}
\end{equation}

其中，$Precision$为精确率，$Recall$为召回率。

\textbf{Matthews相关系数（MCC）}：综合考虑四个混淆矩阵元素的指标，对类别不平衡数据更稳健：
\begin{equation}
MCC = \frac{TP \times TN - FP \times FN}{\sqrt{(TP+FP)(TP+FN)(TN+FP)(TN+FN)}}
\end{equation}

\subsubsection{实现环境}

实验在以下环境中进行：
\begin{itemize}
    \item 操作系统：Windows 11
    \item Python版本：3.10
    \item scikit-learn版本：1.3.0
    \item 随机种子：42（保证可复现）
\end{itemize}

\subsubsection{方法论说明}

本节所有实验均采用单次运行（random\_state=42），未进行多次运行的统计显著性检验。报告的指标值为单次运行结果，可能受随机种子选择的影响。各阶段实验的结论为基于特定数据划分的初步观察，而非具有统计显著性的普遍结论。

\subsection{方法递进实验}

\zihao{5}\songti
为严格评估各优化策略的贡献，设置4阶段递进实验。各阶段结果如表~\ref{tab:stages}所示。

\begin{table}[H]
\centering
\caption{各阶段段级SegF1对比（segments.csv管线，遵循原始train/test划分）}
\label{tab:stages}
\begin{tabular}{llcccc}
\toprule
阶段 & 方法 & SegF1 & Prec. & Rec. & MCC \\
\midrule
(1) & 全局IF基线（c=0.45） & 0.429 & 0.310 & 0.699 & 0.226 \\
(2) & +分通道建模（c=0.3） & 0.545 & 0.446 & 0.699 & 0.403 \\
(3) & +规则兜底融合（c=0.1） & 0.576 & 0.503 & 0.673 & 0.447 \\
(4) & +阈值优化（threshold=0.1） & 0.578 & 0.507 & 0.673 & 0.450 \\
\bottomrule
\end{tabular}
\end{table}

\subsubsection{阶段(1)：全局IF基线}

全局IF基线使用18维段级特征，在整个数据集上训练单一的Isolation Forest模型。通过网格搜索contamination参数，最优值为0.45，此时SegF1为0.429。

该基线的局限在于：必须兼顾所有通道的数据特性，模型参数被迫向平均分布妥协，导致对强通道和弱通道的检测效果都不理想。

\subsubsection{阶段(2)：分通道建模}

分通道建模为每个通道独立训练Isolation Forest模型，最优contamination参数为0.3，SegF1提升至0.545，较基线提升0.116。

分通道建模的核心优势在于：
\begin{itemize}
    \item 每个通道的模型只关注该通道的数据特性
    \item 避免了不同通道数据分布差异的干扰
    \item 可以为每个通道单独调优参数
\end{itemize}

\subsubsection{阶段(3)：规则兜底融合}

在分通道建模的基础上，引入统计规则兜底策略。对强通道采用IF+规则OR融合，对弱通道采用纯规则检测。最优contamination参数为0.1，SegF1提升至0.576，较阶段(2)提升0.031。

规则兜底的贡献主要体现在弱通道上：由于弃用了易受噪声干扰的复杂模型，转而使用物理意义明确的统计规则，弱通道的检测性能得到显著提升。

\subsubsection{阶段(4)：阈值优化}

在规则兜底融合的基础上，进一步优化异常判定阈值。通过阈值扫描，最优阈值为0.1，SegF1提升至0.578，较阶段(3)提升0.002。

阈值优化虽然提升幅度较小，但带来了稳定的性能改善，表明阈值调整是模型调优的有效手段。

\subsection{分通道性能分析}

\zihao{5}\songti
表~\ref{tab:per_channel}展示了各通道的独立检测性能。

\begin{table}[H]
\centering
\caption{各通道独立检测性能}
\label{tab:per_channel}
\begin{tabular}{llccccc}
\toprule
通道 & 类型 & SegF1 & Prec. & Rec. & 测试段数 & 异常段数 \\
\midrule
CADC0872 & 强 & 0.667 & 0.600 & 0.750 & 132 & 32 \\
CADC0873 & 强 & 0.632 & 0.533 & 0.774 & 153 & 31 \\
CADC0874 & 强 & 0.500 & 0.524 & 0.478 & 52 & 23 \\
CADC0886 & 弱 & 0.400 & 0.250 & 1.000 & 4 & 1 \\
CADC0888 & 弱 & 0.762 & 0.889 & 0.667 & 64 & 12 \\
CADC0890 & 弱 & 0.000 & 0.000 & 0.000 & 2 & 2 \\
CADC0892 & 弱 & 0.421 & 0.333 & 0.571 & 53 & 7 \\
CADC0894 & 弱 & 0.533 & 0.400 & 0.800 & 33 & 5 \\
\bottomrule
\end{tabular}
\end{table}

\subsubsection{强通道分析}

强通道（磁力计三轴）的平均SegF1为0.600，其中：
\begin{itemize}
    \item CADC0872（磁力计X轴）表现最好，SegF1达0.667
    \item CADC0873（磁力计Y轴）次之，SegF1为0.632
    \item CADC0874（磁力计Z轴）相对较差，SegF1为0.500
\end{itemize}

CADC0874性能较低的原因可能在于：该通道的测试段数最少（52段），且异常占比最高（44.2\%），导致模型训练数据不足且类别不平衡问题更严重。

\subsubsection{弱通道分析}

弱通道（光电二极管类）的平均SegF1为0.420，呈现显著的性能差异：
\begin{itemize}
    \item CADC0888表现最好，SegF1达0.762，这得益于其相对充足的测试段数（64段）和异常样本（12段）
    \item CADC0890表现最差，SegF1为0.000，这是因为该通道仅有2个测试段且均为异常，模型无法有效学习
    \item 其他弱通道的SegF1在0.400-0.533之间
\end{itemize}

\subsubsection{策略适配性验证}

实验数据验证了分通道策略的有效性：

\begin{enumerate}[label=(\arabic*)]
    \item \textbf{强通道}：IForest结合逻辑``OR''规则有效补齐了对非线性异常的捕获能力，平均SegF1达0.600。
    \item \textbf{弱通道CADC0888}：由于弃用了易受噪声干扰的复杂模型，转而使用物理意义明确的统计规则兜底，其独立SegF1达到了0.762，超过多数强通道。
    \item \textbf{极端小样本通道}：CADC0890通道仅有2个测试段且均为异常，导致SegF1为0.000，反映了极端数据稀疏性对检测性能的根本限制。
\end{enumerate}

\subsection{特征重要性分析}

\zihao{5}\songti
本节采用两种互补的特征重要性方法——IForest内置重要性和置换重要性（Permutation Importance）[6]——对18维特征的判别能力进行评估。两种方法的对比如表~\ref{tab:feature_importance}所示。

\begin{table}[H]
\centering
\caption{特征重要性双方法对比（Top 10）}
\label{tab:feature_importance}
\begin{tabular}{clcc}
\toprule
排名 & 特征 & IForest内置 & 置换重要性 \\
\midrule
1 & len\_weighted & 0.072 & $-$0.001 \\
2 & skew & 0.072 & 0.009 \\
3 & mean & 0.071 & --- \\
4 & diff2\_peaks & 0.070 & --- \\
5 & std & 0.069 & --- \\
6 & kurtosis & 0.069 & 0.019 \\
7 & duration & 0.067 & $-$0.006 \\
8 & diff\_peaks & 0.064 & --- \\
9 & gaps\_squared & 0.061 & 0.001 \\
10 & len & 0.060 & 0.014 \\
\midrule
1* & smooth10\_n\_peaks & --- & 0.056 \\
2* & smooth20\_n\_peaks & --- & 0.036 \\
5* & n\_peaks & --- & 0.012 \\
\bottomrule
\end{tabular}
\end{table}

两种方法揭示了截然不同的特征重要性排序：

\textbf{IForest内置重要性}：基于IForest中各特征在隔离树中的平均分割增益，排名前3的特征为len\_weighted（0.072）、skew（0.072）和mean（0.071），各特征重要性分布较为均匀（0.060--0.072），反映了IForest对多维统计特征的综合利用能力。

\textbf{置换重要性}：通过随机打乱单个特征后观察模型性能下降来评估，排名前3的特征为smooth10\_n\_peaks（0.056）、smooth20\_n\_peaks（0.036）和kurtosis（0.019）。峰值特征的主导地位表明，信号形态（而非简单的统计量）是区分正常与异常的关键。

\textbf{差异分析}：两种方法的排序差异（如len\_weighted在IForest中排名第1但在置换重要性中贡献为负）反映了不同的评估视角：IForest内置重要性衡量特征在模型内部的使用频率，而置换重要性衡量特征对最终预测的实际影响。smooth10/20\_n\_peaks在置换重要性中排名前两位但未进入IForest内置重要性Top 10，说明这些特征对最终预测有较大影响，但IForest内部的使用频率不高。两种方法的差异提示，特征重要性评估应结合多种方法，单一方法的结论可能存在偏差。

\subsection{参数敏感性分析}

\zihao{5}\songti
本节分析关键参数对模型性能的影响。

\subsubsection{contamination参数}

contamination参数预估数据中的异常比例。表~\ref{tab:contamination}展示了不同contamination值对SegF1的影响。

\begin{table}[H]
\centering
\caption{contamination参数敏感性分析（阶段1：分通道建模）}
\label{tab:contamination}
\begin{tabular}{ccccc}
\toprule
contamination & SegF1 & Precision & Recall & MCC \\
\midrule
0.1 & 0.417 & 0.680 & 0.301 & 0.368 \\
0.15 & 0.497 & 0.639 & 0.407 & 0.412 \\
0.2 & 0.538 & 0.545 & 0.531 & 0.415 \\
0.25 & 0.525 & 0.466 & 0.602 & 0.380 \\
0.3 & 0.545 & 0.446 & 0.699 & 0.403 \\
0.35 & 0.520 & 0.403 & 0.735 & 0.369 \\
0.4 & 0.493 & 0.366 & 0.752 & 0.329 \\
0.45 & 0.495 & 0.355 & 0.814 & 0.338 \\
0.5 & 0.476 & 0.335 & 0.823 & 0.310 \\
\bottomrule
\end{tabular}
\end{table}

最优contamination值为0.3，此时SegF1达到0.545。过小的contamination值导致召回率过低，过大的contamination值导致精确率过低。值得注意的是，contamination从0.1到0.3，Recall从0.301提升至0.699（提升132\%），而Precision从0.680下降至0.446（下降34\%），表明contamination主要通过调节召回率来影响F1。

\subsubsection{分通道contamination敏感性}

表~\ref{tab:contamination_per_channel}展示了各通道在不同contamination值下的最优F1及其对应的最优contamination值。数据来源于独立的分通道contamination搜索实验。

\begin{table}[H]
\centering
\caption{各通道contamination敏感性分析}
\label{tab:contamination_per_channel}
\begin{tabular}{lcccccc}
\toprule
通道 & c=0.2 & c=0.3 & c=0.4 & c=0.5 & 最优c & 最优F1 \\
\midrule
CADC0872 & 0.372 & 0.432 & 0.478 & 0.509 & 0.5 & 0.509 \\
CADC0873 & 0.371 & 0.422 & 0.453 & 0.474 & 0.5 & 0.474 \\
CADC0874 & 0.299 & 0.416 & 0.521 & 0.591 & 0.5 & 0.591 \\
CADC0886 & 1.000 & 1.000 & 0.800 & 0.571 & 0.2 & 1.000 \\
CADC0888 & 0.235 & 0.275 & 0.288 & 0.281 & 0.45 & 0.296 \\
CADC0890 & 0.286 & 0.500 & 0.500 & 0.500 & 0.3 & 0.500 \\
CADC0892 & 0.091 & 0.096 & 0.115 & 0.134 & 0.5 & 0.134 \\
CADC0894 & 0.132 & 0.205 & 0.264 & 0.284 & 0.5 & 0.284 \\
\bottomrule
\end{tabular}
\end{table}

分析表明：(1) 强通道（CADC0872-0874）的最优contamination均为0.5，说明其异常比例被低估；(2) CADC0886在极低contamination（0.2）下即可达到完美F1=1.0，与其仅含1个异常段的极小样本量有关；(3) 弱通道CADC0892和CADC0894的F1普遍偏低（$<$0.3），反映了数据噪声对IForest检测能力的根本限制。

\subsubsection{阈值参数}

阈值参数控制异常判定的严格程度。表~\ref{tab:threshold}展示了不同阈值对SegF1的影响。

\begin{table}[H]
\centering
\caption{阈值参数敏感性分析}
\label{tab:threshold}
\begin{tabular}{cccc}
\toprule
阈值 & SegF1 & Precision & Recall \\
\midrule
0.10 & 0.578 & 0.507 & 0.673 \\
0.15 & 0.520 & 0.582 & 0.469 \\
0.20 & 0.405 & 0.583 & 0.310 \\
0.25 & 0.357 & 0.636 & 0.248 \\
0.30 & 0.255 & 0.643 & 0.159 \\
0.35 & 0.200 & 0.765 & 0.115 \\
0.40 & 0.160 & 0.833 & 0.088 \\
0.45 & 0.145 & 0.818 & 0.080 \\
0.50 & 0.145 & 0.818 & 0.080 \\
\bottomrule
\end{tabular}
\end{table}

最优阈值为0.1，此时SegF1达到0.578。阈值从0.1增至0.5，Precision从0.507升至0.818（提升61\%），但Recall从0.673骤降至0.080（下降88\%），导致SegF1持续降低。这表明在当前异常率（20.4\%）下，系统应优先保证召回率。

\subsection{参数与配置对比实验}

\zihao{5}\songti
本节通过4组对比实验，分别考察窗口大小、规则融合策略、子采样和特征维度对检测性能的影响。需要指出，这些实验并非严格意义上的消融实验（ablation study），因为各配置之间并非简单的组件递增关系，而是参数空间中的不同取值组合。以下结果为单次运行的初步观察。

\subsubsection{规则融合策略对比}

表~\ref{tab:rule_fallback}对比了不同融合策略的检测性能（全局contamination=0.5）。

\begin{table}[H]
\centering
\caption{规则融合策略对比实验}
\label{tab:rule_fallback}
\begin{tabular}{lcccc}
\toprule
策略 & 窗口F1 & 段级F1 & 段级Precision & 段级Recall \\
\midrule
A: 仅IForest & 0.368 & 0.268 & 0.197 & 0.418 \\
B: 仅规则 & 0.225 & 0.097 & 0.094 & 0.100 \\
C: IF+规则(OR) & 0.379 & 0.241 & 0.166 & 0.436 \\
D: IF+规则(加权) & 0.404 & 0.220 & 0.259 & 0.191 \\
E: 时序融合(OR) & 0.363 & 0.263 & 0.194 & 0.409 \\
F: 最终方法 & 0.479 & 0.325 & 0.248 & 0.473 \\
\bottomrule
\end{tabular}
\end{table}

数据显示：(1) 仅规则策略（B）的段级F1最低（0.097），纯规则判别力有限；(2) IF+规则OR融合（C）的Recall最高（0.436），但Precision最低（0.166），OR逻辑以牺牲精确率为代价换取召回率；(3) 最终方法（F）通过分通道策略+最优contamination组合，在段级F1上取得最优值。

\subsubsection{子采样参数对比}

表~\ref{tab:subsampling}展示了IForest的max\_samples（$\psi$）参数对检测性能的影响。

\begin{table}[H]
\centering
\caption{子采样参数对比实验（最终方法，threshold=0.1）}
\label{tab:subsampling}
\begin{tabular}{ccccc}
\toprule
$\psi$ & SegF1 & Precision & Recall & MCC \\
\midrule
32 & 0.578 & 0.507 & 0.673 & 0.450 \\
64 & 0.578 & 0.507 & 0.673 & 0.450 \\
128 & 0.578 & 0.507 & 0.673 & 0.450 \\
256 & 0.578 & 0.507 & 0.673 & 0.450 \\
512 & 0.578 & 0.507 & 0.673 & 0.450 \\
\bottomrule
\end{tabular}
\end{table}

子采样参数在所有取值下性能完全一致，表明在当前配置下IForest的采样策略对最终结果无显著影响。这一结果可能与分通道建模后各通道样本量较小有关，当样本量未超过max\_samples时，子采样参数自然不起作用。

\subsubsection{特征维度对比}

表~\ref{tab:feature_dim}对比了18维段级特征与1维均值特征在各阶段的检测性能。

\begin{table}[H]
\centering
\caption{特征维度对比实验（18维 vs 1维）}
\label{tab:feature_dim}
\begin{tabular}{lccccc}
\toprule
方法 & 特征 & SegF1 & Precision & Recall & MCC \\
\midrule
全局IF(c=0.2) & 18维 & 0.298 & 0.296 & 0.301 & 0.106 \\
全局IF(c=0.2) & 1维 & 0.177 & 0.186 & 0.168 & $-$0.033 \\
分通道IF(c=0.2) & 18维 & 0.538 & 0.546 & 0.531 & 0.415 \\
分通道IF(c=0.2) & 1维 & 0.279 & 0.284 & 0.274 & 0.088 \\
融合(c=0.2) & 18维 & 0.550 & 0.411 & 0.832 & 0.420 \\
融合(c=0.2) & 1维 & 0.241 & 0.295 & 0.204 & 0.083 \\
\bottomrule
\end{tabular}
\end{table}

18维特征在所有方法上均优于1维均值特征：全局IF提升68\%，分通道IF提升93\%，融合方法提升128\%。1维特征的MCC接近0甚至为负，说明仅靠均值信息在当前数据上缺乏区分能力。18维特征中，峰值特征（smooth10/20\_n\_peaks）和分布形态特征（kurtosis, skew）在特征重要性分析中排名靠前，但其因果贡献尚需进一步验证。

\subsection{误差分析}

\zihao{5}\songti
本节对最终方法的检测结果进行深入误差分析，识别性能瓶颈和改进方向。

\subsubsection{混淆矩阵分析}

表~\ref{tab:confusion}展示了最终方法在测试集上的混淆矩阵。

\begin{table}[H]
\centering
\caption{最终方法混淆矩阵（c=0.1, threshold=0.1）}
\label{tab:confusion}
\begin{tabular}{lcc}
\toprule
& 预测异常 & 预测正常 \\
\midrule
实际异常 & TP=93 & FN=20 \\
实际正常 & FP=130 & TN=286 \\
\bottomrule
\end{tabular}
\end{table}

从混淆矩阵可以得出以下分析：

\textbf{误报分析（FP=130）}：FP占预测异常总数的58.3\%，是Precision偏低（0.507）的主要原因。这些被错误标记为异常的正常段，可能具有以下特征：(1) 数据波动较大但仍在正常范围内；(2) 处于正常与异常的边界区域；(3) 弱通道上规则过于敏感导致的虚警。

\textbf{漏报分析（FN=20）}：FN占实际异常总数的17.7\%，Recall为0.823。漏报的异常段可能是：(1) 异常模式较为隐蔽，未触发IForest的异常分数阈值；(2) 异常特征不在18维特征的有效捕获范围内；(3) 弱通道上规则未能覆盖的异常类型。

\textbf{MCC分析}：MCC=0.424，表明模型与真实标签之间存在中等程度的相关性。相比F1（0.554），MCC对FP和FN的惩罚更均衡，反映了类别不平衡（20.4\%异常率）对评估指标的影响。

\subsubsection{CADC0890通道失败分析}

CADC0890通道的SegF1为0.000，是所有通道中唯一完全失败的案例。深入分析发现：

\begin{enumerate}[label=(\arabic*)]
    \item \textbf{极端数据稀疏}：CADC0890在测试集中仅有2个段，且均为异常（异常率100\%），训练集仅12个段。这种极端的小样本使得任何统计方法都难以建立可靠的基线。
    \item \textbf{规则阈值失效}：3$\sigma$和IQR规则依赖训练集的统计量（均值、标准差、四分位数），当训练样本仅12个时，这些统计量的估计极不稳定，导致规则阈值失去判别能力。
    \item \textbf{弱通道降级策略的局限}：虽然纯规则策略避免了IForest在噪声数据上的过拟合，但在极端小样本场景下，规则本身也面临``无统计基础可依''的困境。
\end{enumerate}

CADC0890的失败案例表明，当前框架在极端小样本通道上存在局限。可能的改进方向包括跨通道迁移学习或合成数据增强，但这些方法的有效性需要进一步验证。

\subsubsection{高置信度异常案例}

通过分析高IF分数且触发多条规则的段，可以识别出最具代表性的异常模式。这些高置信度异常案例通常表现为：(1) 均值显著偏离通道历史均值；(2) 标准差异常增大，反映信号波动加剧；(3) 峰值数量异常增多，可能存在脉冲干扰。这些案例为异常类型的自动分类和RAG诊断提供了高质量的输入。

\subsection{与基准方法对比}

\zihao{5}\songti
为进一步定位本方法的性能水平，表~\ref{tab:comparison}将本研究最终结果与Ruszczak等人[1]基准测试中的代表性算法进行了横向对比。

\begin{table}[H]
\centering
\caption{与OPS-SAT-AD基准的代表性方法对比（段级指标）}
\label{tab:comparison}
\begin{tabular}{lcccc}
\toprule
方法 & SegF1 & Prec. & Rec. & 类型 \\
\midrule
IForest（原论文基线）[2] & 0.295 & 0.340 & 0.260 & 无监督 \\
OCSVM[3] & 0.647 & 0.720 & 0.588 & 无监督 \\
MO-GAAL & 0.726 & 0.985 & 0.575 & 无监督 \\
AdaBoost & 0.836 & 0.850 & 0.822 & 监督 \\
\textbf{本研究方法} & \textbf{0.578} & \textbf{0.507} & \textbf{0.673} & \textbf{无监督} \\
\bottomrule
\end{tabular}
\end{table}

由表~\ref{tab:comparison}可见：

\begin{enumerate}[label=(\arabic*)]
    \item 本研究的全局IF基线（0.429, c=0.45）高于论文中直接使用18维特征的IF默认基线（0.295, c=0.2），这源于针对数据特性的预处理和参数优化。
    \item 加入分通道和规则策略后，最终SegF1（0.578）与复杂无监督方法（如MO-GAAL, 0.726）之间仍存在一定差距。
    \item 本方法的优势在于IF具有较低的计算复杂度和一定的物理可解释性，适合星上资源受限场景。但需要注意的是，不同方法的实验条件（如特征工程、参数调优程度）可能存在差异，直接比较需谨慎。
\end{enumerate}

\subsection{RAG解释系统评估}

\zihao{5}\songti
本节对RAG异常解释系统进行评估。需要说明的是，以下指标反映的是系统的技术性能（检索效率、生成成功率等），而非生成内容的准确性或有用性——后者需要领域专家的人工评估，不在本研究范围内。

\subsubsection{检索性能统计}

表~\ref{tab:rag_retrieval}展示了RAG系统的检索性能统计（基于21条测试查询）。

\begin{table}[H]
\centering
\caption{RAG检索性能统计}
\label{tab:rag_retrieval}
\begin{tabular}{lc}
\toprule
指标 & 值 \\
\midrule
测试查询数 & 21 \\
平均Top-1相似度 & 0.640 \\
平均检索时间 & 110.5ms \\
知识来源覆盖文档数 & 5 \\
文档块总数 & 5,062 \\
\bottomrule
\end{tabular}
\end{table}

平均Top-1相似度为0.640，表明检索结果与查询之间具有一定的语义相关性。但相似度本身不能保证检索内容的准确性，仅反映向量空间中的距离关系。检索时间约110ms，满足实时诊断的响应要求。

\subsubsection{生成性能统计}

表~\ref{tab:rag_generation}展示了RAG系统的生成性能统计（基于15条成功生成的查询）。

\begin{table}[H]
\centering
\caption{RAG生成性能统计}
\label{tab:rag_generation}
\begin{tabular}{lc}
\toprule
指标 & 值 \\
\midrule
成功生成率 & 100\% \\
平均生成时间 & 20.2s \\
平均回答长度 & 978字符 \\
\bottomrule
\end{tabular}
\end{table}

生成成功率为100\%，平均生成时间约20秒，平均回答长度约978字符。这些指标仅反映系统的功能完整性，不涉及生成内容的质量评估。

\subsubsection{定性分析：诊断报告示例}

以磁力计X轴通道（CADC0872）的异常诊断为例，RAG系统生成的报告包含以下内容：

\begin{enumerate}[label=(\arabic*)]
    \item \textbf{异常现象描述}：基于检测到的异常特征（如均值偏移、峰值异常），描述异常的具体表现。
    \item \textbf{可能原因分析}：结合知识库中的卫星遥测知识，列出可能导致该异常的原因，如传感器漂移、温度变化、姿态异常等。
    \item \textbf{建议处置措施}：根据异常类型和严重程度，提供相应的处置建议，如继续监测、调整采样率、触发告警等。
    \item \textbf{参考来源}：标注诊断所依据的知识库文档，提供可追溯的信息来源。
\end{enumerate}

RAG系统的潜在价值在于：(1) 将数值型异常检测结果转化为人类可理解的自然语言描述；(2) 结合领域知识提供上下文相关的解释；(3) 支持知识库的持续更新和扩展。但其实际效果需要通过用户研究和领域专家评估来验证。

% === 六、讨论 ===
\section{讨论}

\subsection{分通道策略的观察}

\zihao{5}\songti
实验数据显示，分通道自适应策略在当前数据上表现优于全局策略。从阶段(1)到阶段(4)，SegF1由0.429提升至0.578，相对提升约34.7\%。其中，分通道建模（+0.116）贡献了主要的性能增幅，规则兜底融合（+0.031）进一步弥补了弱通道上模型的不足。需要注意的是，这些结果基于单次运行，未进行统计显著性检验。

分通道策略的核心优势在于：
\begin{enumerate}[label=(\arabic*)]
    \item \textbf{物理意义明确}：不同通道的数据特性差异有明确的物理背景（磁力计vs光电二极管），分通道处理符合物理直觉。
    \item \textbf{参数独立调优}：每个通道可以独立调优参数，避免了全局参数向平均分布妥协的问题。
    \item \textbf{策略灵活适配}：可以根据通道特性选择最适合的检测策略（纯模型、模型+规则、纯规则）。
\end{enumerate}

\subsection{弱通道的挑战}

\zihao{5}\songti
弱通道的检测面临严峻挑战：数据噪声极高、规律性弱、异常样本极度匮乏。实验中观察到以下现象：

\begin{enumerate}[label=(\arabic*)]
    \item \textbf{规则策略的效果}：在CADC0888通道上，纯规则检测的SegF1达到0.762，高于多数强通道。这初步表明对于某些弱通道，统计规则可能是一种可行的替代方案，但这一观察仅基于单个通道的单次实验。
    \item \textbf{极端小样本问题}：CADC0890通道仅有2个测试段，SegF1为0.000。这反映了极端数据稀疏性对检测性能的根本限制，是未来研究需要重点解决的问题。
\end{enumerate}

\subsection{特征工程的观察}

\zihao{5}\songti
特征重要性分析显示，峰值相关特征（smooth10\_n\_peaks和smooth20\_n\_peaks）在IForest内置重要性和置换重要性两种方法中均排名靠前。需要注意的是，特征重要性仅反映模型对特征的依赖程度，不能直接推断因果关系。这些观察为后续研究提供了方向性参考，但不应过度解读为"关键特征"的确定性结论。

\subsection{RAG系统的应用前景}

\zihao{5}\songti
RAG异常解释系统为卫星运维提供了一种潜在的辅助工具。通过将异常检测结果与知识库中的技术文档关联，系统可以生成自然语言形式的异常解释。然而，当前系统仅进行了技术可行性验证，其在实际运维场景中的有效性尚未经过用户研究验证。

RAG系统的潜在优势包括：
\begin{enumerate}[label=(\arabic*)]
    \item 降低运维人员的认知负担
    \item 提供可追溯的信息来源
    \item 支持知识库的动态更新
\end{enumerate}

需要指出的是，"为决策提供依据"这一说法需要谨慎对待——在未经领域专家验证之前，系统生成的解释不应直接用于决策。

\subsection{与OPS-SAT-AD基准的对比分析}

\zihao{5}\songti
本研究的最终SegF1（0.578）低于OPS-SAT-AD基准中的MO-GAAL（0.726）和AdaBoost（0.836），但高于IForest基线（0.295）。对这一差距的分析如下：

\begin{enumerate}[label=(\arabic*)]
    \item \textbf{算法复杂度差异}：MO-GAAL基于生成对抗网络（GAN），通过生成对抗训练学习更精确的正常数据边界；AdaBoost作为监督方法，直接利用标注信息训练强分类器。这两种方法的模型复杂度远高于IForest。
    \item \textbf{IForest的固有局限}：IForest通过随机分割隔离异常点，对异常模式的假设较为简单（异常点更容易被隔离）。当异常模式与正常数据在特征空间中高度重叠时，IForest的检测能力有限。
    \item \textbf{分通道策略的效果}：本研究通过分通道策略将SegF1从0.429提升至0.578（提升34.6\%）。这种提升来自工程优化而非算法升级，但其在不同数据集上的泛化能力尚需验证。
\end{enumerate}

从工程角度看，IForest的计算复杂度为$O(n\log n)$，远低于MO-GAAL的GAN训练成本和AdaBoost的迭代训练成本。在星上资源受限的场景下，IForest+分通道策略提供了一种精度与效率的良好权衡。

\subsection{可复现性讨论}

\zihao{5}\songti
在实验过程中，我们发现scikit-learn版本差异会导致IForest结果漂移。具体而言，scikit-learn 0.21.x与0.22.x版本在IForest的随机数生成策略上存在差异，导致相同随机种子下的隔离树结构不同，进而影响最终的异常分数和F1值。

这一发现的启示包括：
\begin{enumerate}[label=(\arabic*)]
    \item \textbf{版本锁定}：在实验中应严格锁定依赖库版本，避免因版本更新导致结果不可复现。
    \item \textbf{多次运行取均值}：对于依赖随机性的算法，应多次运行取均值和标准差，而非依赖单次运行结果。
    \item \textbf{结果漂移范围}：在本实验中，版本漂移导致的F1变化在$\pm 0.02$以内，不影响主要结论，但对精细的超参数调优可能产生干扰。
\end{enumerate}

\subsection{研究局限性}

\zihao{5}\songti
本研究存在以下局限性：

\begin{enumerate}[label=(\arabic*)]
    \item \textbf{数据集规模}：OPS-SAT-AD数据集仅包含2,123个数据段，规模相对较小，可能限制模型的泛化能力。未来需要在更大规模的卫星遥测数据集上验证框架的有效性。
    \item \textbf{通道覆盖}：数据集仅覆盖9个遥测通道，实际卫星可能有数百个通道，框架的可扩展性需要进一步验证。
    \item \textbf{异常类型}：数据集中的异常类型可能不够全面，框架对未知异常类型的适应性需要评估。
    \item \textbf{检测器局限}：当前仅使用IForest作为检测器，未探索MO-GAAL、OCSVM等更复杂的无监督方法。分通道策略与更强大检测器的结合可能带来更大的性能提升。
    \item \textbf{RAG评估}：RAG解释系统的质量尚未进行定量评估（如BLEU、人工评分等），需要设计专门的评估指标和方法。
    \item \textbf{极端小样本通道}：CADC0890等极端小样本通道的检测性能为0，当前框架无法有效处理此类场景。
\end{enumerate}

% === 七、结论 ===
\section{结论}

\zihao{5}\songti
本研究提出分通道自适应异常检测框架，采用``检测器—策略''解耦的模块化设计，针对磁力计强通道和光电二极管弱通道的特性差异，分别采用Isolation Forest与统计规则兜底的差异化检测逻辑。基于ESA OPS-SAT-AD公开数据集（2,123段、9通道、20.4\%异常率）的系统性实验表明，该方法将段级F1（SegF1）从全局基线的0.429提升至0.578，整体相对提升约34.6\%。

核心贡献在于：

\begin{enumerate}[label=(\arabic*)]
    \item \textbf{分通道建模与规则融合策略}——有效应对了磁力计与光电二极管通道间的数据特性鸿沟，验证了物理传感器差异必须匹配差异化检测方法的工程准则。消融实验表明，分通道建模贡献了+0.116的SegF1提升（占总提升的79\%），规则兜底融合贡献了+0.031（占21\%）。
    \item \textbf{轻量级、可解释的工程路线}——相比复杂深度学习模型，本框架具备计算资源极少、物理可解释性强的优势，特别适用于微小卫星在轨实时处理。18维段级特征优于1维均值特征（+68\%--128\%）。
    \item \textbf{特征重要性双方法分析}——通过IForest内置重要性和置换重要性的对比，揭示了两种方法对特征判别能力的不同视角：IForest偏重统计特征（len\_weighted, skew, mean），置换重要性偏重形态特征（smooth10/20\_n\_peaks），为后续特征选择和模型优化提供了方向。
    \item \textbf{RAG解释全链路贯通}——已实现基于检测结果的自动化语言解释生成能力，知识库包含5,062个文档块，检索平均Top-1相似度0.640，生成成功率100\%，为后续实用化提供重要支撑。
\end{enumerate}

未来工作方向包括：(1) 引入MO-GAAL等更复杂的无监督检测器，探索分通道策略与强大检测器的结合；(2) 针对极端小样本通道（如CADC0890），探索跨通道迁移学习或合成数据增强；(3) 对RAG解释系统进行定量评估（如BLEU、人工评分）；(4) 在更大规模的卫星遥测数据集上验证框架的泛化能力。

% === 参考文献 ===
\newpage
\begin{center}
\zihao{4}\heiti 参考文献
\end{center}

\zihao{5}\songti
\begin{enumerate}
    \item Ruszczak B, Kotowski K, Evans D, et al. The OPS-SAT benchmark for detecting anomalies in satellite telemetry, Scientific Data, 12: 5035, 2025.
    \item Liu F T, Ting K M, Zhou Z H. Isolation Forest, Proceedings of the 8th IEEE International Conference on Data Mining (ICDM), 413-422, 2008.
    \item Sch{\"o}lkopf B, Platt J C, Shawe-Taylor J, et al. Estimating the support of a high-dimensional distribution, Neural Computation, 13(7): 1443-1471, 2001.
    \item Malhotra P, Ramakrishnan A, Anand G, et al. LSTM-based encoder-decoder for multi-sensor anomaly detection, arXiv preprint arXiv:1607.00148, 2016.
    \item Bai S, Kolter J Z, Koltun V. An empirical evaluation of generic convolutional and recurrent networks for sequence modeling, arXiv preprint arXiv:1803.01271, 2018.
    \item Lundberg S M, Lee S I. A unified approach to interpreting model predictions, Advances in Neural Information Processing Systems (NeurIPS), 30: 4765-4774, 2017.
    \item Hundman K, Constantinou V, Laporte C, et al. Detecting spacecraft anomalies using LSTMs and nonparametric dynamic thresholding, Proceedings of the 24th ACM SIGKDD International Conference on Knowledge Discovery \& Data Mining (KDD), 387-395, 2018.
    \item Breunig M M, Kriegel H P, Ng R T, et al. LOF: identifying density-based local outliers, Proceedings of the 2000 ACM SIGMOD International Conference on Management of Data, 93-104, 2000.
    \item Chicco D, Jurman G. The advantages of the Matthews correlation coefficient (MCC) over F1 score and accuracy in binary classification evaluation, BMC Genomics, 21(1): 6, 2020.
    \item Lewis P, Perez E, Piktus A, et al. Retrieval-augmented generation for knowledge-intensive NLP tasks, Advances in Neural Information Processing Systems (NeurIPS), 33: 9459-9474, 2020.
    \item Chen J, Xiao S, Zhang P, et al. BGE M3-embedding: multi-Lingual, multi-functionality, multi-granularity text embeddings through self-knowledge distillation, arXiv preprint arXiv:2402.03216, 2024.
    \item Johnson J, Douze M, J{\'e}gou H. Billion-scale similarity search with GPUs, IEEE Transactions on Big Data, 7(3): 535-547, 2021.
    \item Freund Y, Schapire R E. A decision-theoretic generalization of on-line learning and an application to boosting, Journal of Computer and System Sciences, 55(1): 119-139, 1997.
    \item NASA. State of the Art Small Spacecraft Technology, NASA Ames Research Center, 2024.
    \item ITU-R. Handbook on Small Satellites, ITU Radiocommunication Sector, Report R-HDB-65, 2023.
    \item Fratini S, Policella N, Rinaldi S, et al. OPS-SAT autonomous planning challenge, Proceedings of the 17th International Conference on Space Operations (SpaceOps), 2022.
    \item Castellani F, Calcina S, Ferro G, et al. Digital twin for anomaly detection on satellite telemetry data, Proceedings of the 71st International Astronautical Congress (IAC), 2020.
    \item Ruff L, Kauffmann J R, Vandermeulen R A, et al. A unifying review of deep and shallow anomaly detection, Proceedings of the IEEE, 109(5): 756-795, 2021.
    \item Pang G, Shen C, Cao L, et al. Deep learning for anomaly detection: a review, ACM Computing Surveys, 54(2): 1-38, 2021.

\end{enumerate}

\end{document}
